\section{Experiments}\label{sec:exp}

In all experiments the trained optimizers use two-layer LSTMs with 20 hidden
units in each layer. Each optimizer is trained by minimizing
Equation~\ref{eq:meta-optimization} using truncated BPTT as described in
Section~\ref{sec:learning-to-learn}. The minimization is performed using ADAM
with a learning rate chosen by random search.

We use early stopping when training the optimizer in order to avoid overfitting
the optimizer.  After each epoch (some fixed number of learning steps) we
freeze the optimizer parameters and evaluate its performance.  We pick the best
optimizer (according to the final validation loss) and report its average
performance on a number of freshly sampled test problems.

We compare our trained optimizers with standard optimizers used in Deep
Learning: SGD, RMSprop, ADAM, and Nesterov's accelerated gradient (NAG).
For each of these optimizer and each problem we tuned the learning rate, and
report results with the rate that gives the best final error for each problem.
When an optimizer has more parameters than just a learning rate (e.g. decay
coefficients for ADAM) we use the default values from the \texttt{optim}
package in Torch7.  Initial values of all optimizee parameters were sampled
from an IID Gaussian distribution.



\subsection{Quadratic functions}

In this experiment we consider training an optimizer on a simple class of
synthetic 10-dimensional quadratic functions.  In particular we consider minimizing
functions of the form
\begin{align*}
    f(\theta) = \|W\theta - y\|_2^2
\end{align*}
for different 10x10 matrices $W$ and 10-dimensional vectors $y$ whose elements are drawn
from an IID Gaussian distribution.
Optimizers were trained by optimizing random functions from this family and
tested on newly sampled functions from the same distribution.  Each function was
optimized for 100 steps and the trained optimizers were unrolled for 20 steps.
We have not used any preprocessing, nor postprocessing.

Learning curves for different optimizers, averaged over many functions, are
shown in the left plot of Figure~\ref{fig:mnist-1}.  Each curve corresponds to
the average performance of one optimization algorithm on many test functions;
the solid curve shows the learned optimizer performance and dashed curves show
the performance of the standard baseline optimizers.  It is clear the learned
optimizers substantially outperform the baselines in this setting.

\begin{figure}
\centering
\includegraphics[width=\linewidth]{plots/mnist-1.pdf}
\caption{Comparisons between learned and hand-crafted optimizers performance.
Learned optimizers are shown with solid lines and hand-crafted optimizers are
shown with dashed lines.  Units for the $y$ axis in the MNIST plots are logits.
\textbf{Left:} Performance of different optimizers on randomly sampled
10-dimensional quadratic functions. \textbf{Center:} the LSTM optimizer
outperforms standard methods training the base network on MNIST.
\textbf{Right:} Learning curves for steps 100-200 by an optimizer trained to
optimize for 100 steps (continuation of center plot).}
\label{fig:mnist-1}
\end{figure}

\begin{figure}
\centering
\includegraphics[width=\linewidth]{plots/mnist-2.pdf}
\caption{Comparisons between learned and hand-crafted optimizers performance.
Units for the $y$ axis are logits. \textbf{Left:} Generalization to the
different number of hidden units (40 instead of 20).  \textbf{Center:}
Generalization to the different number of hidden layers (2 instead of 1).  This
optimization problem is very hard, because the hidden layers are very narrow.
\textbf{Right:} Training curves for an MLP with 20 hidden units using ReLU
activations.  The LSTM optimizer was trained on an MLP with sigmoid
activations.}
\label{fig:mnist-2}
\end{figure}

\begin{figure}
\centering
\includegraphics[width=\linewidth]{plots/mnist-sigmoid.pdf}
\caption{Systematic study of final MNIST performance as the optimizee
architecture is varied, using sigmoid non-linearities. The vertical dashed line
in the left-most plot denotes the architecture at which the LSTM is trained and
the horizontal line shows the final performance of the trained optimizer in
this setting.}
\label{fig:mnist-sigmoid}
\end{figure}

\subsection{Training a small neural network on MNIST}

In this experiment we test whether trainable optimizers can learn to optimize a
small neural network on MNIST, and also explore how the trained optimizers
generalize to functions beyond those they were trained on.  To this end, we
train the optimizer to optimize a base network and explore a series of
modifications to the network architecture and training procedure at test time.

%\begin{wrapfigure}[25]{r}{0.6\textwidth}
%\centering
%\includegraphics[width=\linewidth]{plots/mnist-all.pdf}
%\caption{Systematic study of the final MNIST performance as the optimizee
%architecture is varied, using both sigmoid and ReLU non-linearities. Note that
%we train two optimizers, both on MNIST problems using architectures of 1 layer
%and 20 hidden units, but with and without ReLU.}
%\label{fig:mnist-all}
%\end{wrapfigure}

In this setting the objective function $f(\theta)$ is the cross entropy of a
small MLP with parameters $\theta$. The values of $f$ as well as the gradients
$\partial f(\theta) / \partial\theta$ are estimated using random minibatches of
128 examples. The base network is an MLP with one hidden layer of 20 units
using a sigmoid activation function. The only source of variability between
different runs is the initial value $\theta_0$ and randomness in minibatch
selection. Each optimization was run for 100 steps and the trained optimizers
were unrolled for 20 steps. We used input preprocessing described in
Appendix~\ref{sec:pre} and rescaled the outputs of the LSTM by the factor
$0.1$.

Learning curves for the base network using different optimizers are displayed
in the center plot of Figure~\ref{fig:mnist-1}. In this experiment NAG, ADAM,
and RMSprop exhibit roughly equivalent performance the LSTM optimizer
outperforms them by a significant margin. The right plot in
Figure~\ref{fig:mnist-1} compares the performance of the LSTM optimizer if it
is allowed to run for 200 steps, despite having been trained to optimize for
100 steps.  In this comparison we re-used the LSTM optimizer from the previous
experiment, and here we see that the LSTM optimizer continues to outperform the
baseline optimizers on this task.

\paragraph{Generalization to different architectures} Figure~\ref{fig:mnist-2}
shows three examples of applying the LSTM optimizer to train networks with
different architectures than the base network on which it was trained.  The
modifications are (from left to right) (1) an MLP with 40 hidden units instead
of 20, (2) a network with two hidden layers instead of one, and (3) a network
using ReLU activations instead of sigmoid.  In the first two cases the LSTM
optimizer generalizes well, and continues to outperform the hand-designed
baselines despite operating outside of its training regime.  However, changing
the activation function to ReLU makes the dynamics of the learning procedure
sufficiently different that the learned optimizer is no longer able to
generalize. Finally, in Figure~\ref{fig:mnist-sigmoid} we show the results of
systematically varying the tested architecture; for the LSTM results we again
used the optimizer trained using 1 layer of 20 units and sigmoid
non-linearities. Note that in this setting where the test-set problems are
similar enough to those in the training set we see even better generalization
than the baseline optimizers.

\subsection{Training a convolutional network on CIFAR-10}

Next we test the performance of the trained neural optimizers on optimizing
classification performance for the CIFAR-10 dataset \citep{krizhevsky:2009}. In
these experiments we used a model with both convolutional and feed-forward
layers. In particular, the model used for these experiments includes three
convolutional layers with max pooling followed by a fully-connected layer with
32 hidden units; all non-linearities were ReLU activations with batch
normalization.

The coordinatewise network decomposition introduced in
Section~\ref{sec:coordinatewise}---and used in the previous
experiment---utilizes a single LSTM architecture with shared weights, but
separate hidden states, for each optimizee parameter. We found that this
decomposition was not sufficient for the model architecture introduced in this
section due to the differences between the fully connected and convolutional
layers. Instead we modify the optimizer by introducing two LSTMs: one proposes
parameter updates for the fully connected layers and the other updates the
convolutional layer parameters. Like the previous LSTM optimizer we still
utilize a coordinatewise decomposition with shared weights and individual
hidden states, however LSTM weights are now shared only between parameters of
the same type (i.e.\ fully-connected vs.\ convolutional).

The performance of this trained optimizer compared against the baseline
techniques is shown in Figure~\ref{fig:cifar}. The left-most plot displays the
results of using the optimizer to fit a classifier on a held-out test set. The
additional two plots on the right display the performance of the trained
optimizer on modified datasets which only contain a subset of the labels, i.e.
the CIFAR-2 dataset only contains data corresponding to 2 of the 10 labels.
Additionally we include an optimizer \emph{LSTM-sub} which was only trained on
the held-out labels.

In all these examples we can see that the LSTM optimizer learns much more
quickly than the baseline optimizers, with significant boosts in performance
for the CIFAR-5 and especially CIFAR-2 datsets. We also see that the optimizers
trained only on a disjoint subset of the data is hardly effected by this
difference and transfers well to the additional dataset.

\begin{figure}
\centering
\includegraphics[width=\linewidth]{plots/cifar.pdf}
\caption{Optimization performance on the CIFAR-10 dataset and subsets. Shown on
the left is the LSTM optimizer versus various baselines trained on CIFAR-10 and
tested on a held-out test set. The two plots on the right are the performance
of these optimizers on subsets of the CIFAR labels. The additional optimizer
\emph{LSTM-sub} has been trained only on the heldout labels and is hence
transferring to a completely novel dataset.}
\label{fig:cifar}
\end{figure}

\subsection{Neural Art}

\begin{figure}
\centering
\includegraphics[width=0.85\linewidth]{plots/neural-art.pdf}
\vspace{-0.1cm}
\caption{
Optimization curves for Neural Art. Content images come from the test set,
which was not used during the LSTM optimizer training. Note: the y-axis is in
log scale and we zoom in on the interesting portion of this plot.
\textbf{Left:} Applying the training style at the training resolution.
\textbf{Right:} Applying the test style at double the training resolution.}
\label{fig:neural-art}
\end{figure}

\begin{figure}
\centering
\includegraphics[width=0.46\linewidth]{plots/neural-art-img1-big-boat.png}
\qquad
\includegraphics[width=0.46\linewidth]{plots/neural-art-img2-big-scream.png}
\caption{
Examples of images styled using the LSTM optimizer.  Each triple consists of
the content image (left), style (right) and image generated by the LSTM
optimizer (center).  \textbf{Left:} The result of applying the training style
at the training resolution to a test image.  \textbf{Right:} The result of
applying a new style to a test image at double the resolution on which the
optimizer was trained.}
\label{fig:neural-art-examples}
\end{figure}


The recent work on artistic style transfer using convolutional networks, or
Neural Art \citep{gatys:2015}, gives a natural testbed for our method, since
each content and style image pair gives rise to a different optimization
problem.  Each Neural Art problem starts from a \emph{content image}, $c$, and
a \emph{style image}, $s$, and is given by
\begin{align*}
f(\theta) = \alpha \mathcal{L}_\mathrm{content}(c, \theta)
+ \beta \mathcal{L}_\mathrm{style}(s, \theta) + \gamma \calL_{\mathrm{reg}}(\theta)
\,
\end{align*}
The minimizer of $f$ is the \emph{styled image}.
The first two terms try to match the content and style of the styled image to that of their first argument, and the third term is a regularizer that encourages smoothness in the styled image.
Details can be found in \citep{gatys:2015}.

% Neural
% Art is formulated as the following optimization problem.
% Starting from a \emph{content image}, $c$, and a \emph{style image}, $s$, a \emph{styled
% image}, $\theta$, is produced by optimizing the following function
% \begin{align*}
% f(\theta) = \alpha \mathcal{L}_\mathrm{content}(c, \theta)
% + \beta \mathcal{L}_\mathrm{style}(s, \theta) + \gamma ||\theta||_{\mathrm{tv},p}
% \,,
% \end{align*}
% where the content loss $\mathcal{L}_\mathrm{content}(c, \theta)$ measures the distance between the contents of images $c$ and $\theta$,
% $\mathcal{L}_\mathrm{style}(s, \theta)$ measures the distance between the styles of $s$ and $\theta$ and $||\theta||_{\mathrm{tv},p}$ is
% some regularization term.
% The values $\mathcal{L}_\mathrm{content}(c, \theta)$ and $\mathcal{L}_\mathrm{style}(s, \theta)$ are in fact defined using activations in some fixed convolutional neural network when
% applied to appropriate images.
% The initial value of the styled images is set to the content image $\theta_0 = c$.
% The neural network used is the VGG-19 network
% of \citep{simonyan:2014}, which seems to give the most visually appealing
% results.

We train optimizers using only 1 style and 1800 content images taken from
ImageNet \citep{deng:2009}.  We randomly select 100 content images for testing
and 20 content images for validation of trained optimizers.  We train the
optimizer on 64x64 content images from ImageNet and one fixed style image.  We
then test how well it generalizes to a different style image and higher
resolution (128x128).  Each image was optimized for 128 steps and trained
optimizers were unrolled for 32 steps.  Figure~\ref{fig:neural-art-examples}
shows the result of styling two different images using the LSTM optimizer.  The
LSTM optimizer uses inputs preprocessing described in Appendix~\ref{sec:pre}
and no postprocessing. See Appendix~\ref{sec:neural-art-images} for additional images.

Figure~\ref{fig:neural-art} compares the performance of the LSTM optimizer to
standard optimization algorithms.  The LSTM optimizer outperforms all standard
optimizers if the resolution and style image are the same as the ones on which
it was trained.  Moreover, it continues to perform very well when both the
resolution and style are changed at test time.

Finally, in Appendix~\ref{sec:proposed-updates} we qualitatively examine the
behavior of the step directions generated by the learned optimizer.

%We also experimented with changing only one of the resolution or style with similar results.

%testing this trained procedure on (top left) problems from the training set,
%(top right) problems with different styles, (bottom left) problems with
%different target sizes, and (bottom right) problems with both style and size
%differences.
%In all cases we see that the learned optimizer outperforms or at least matches
%the performance of all other methods.
%and is able to generalize beyond the strict style and size
%settings it was trained with. Figure~\ref{fig:neural-art-examples} shows neural art styles using the trained optimizer for the larger target instances.

